% Programming with InitechBase, Chapter 17: Working with Characters (AN-16).
% Panels replayed (examples/p17-*.dot); sources core/fn_builtins.c (LEFT,
% SUBSTR, LOWER, UPPER, TRIM, AT, STUFF), core/eval.c ($), the program
% NAMEFMT.PRG.
\chapter{Working with Characters}
\Gates{test-pubs-examples, test-xbase-fn-a, test-xbase-fn-c}
The staff roster keeps every name in capitals, as the clerks key them.  A
letter to an employee prints the name with only its first letter in capitals.
This chapter builds the salutation of a letter from the fields of the roster.

\begin{ChapterCovers}
\item Taking strings apart: \fn{LEFT()} and \fn{SUBSTR()}
\item Changing case: \fn{UPPER()} and \fn{LOWER()}
\item Searching: \texttt{\$} and \fn{AT()}
\item Replacing part of a string: \fn{STUFF()}
\end{ChapterCovers}

\ATEHead{Preparing for This Chapter}
Start InitechBase with a fresh copy of the sample diskette, and type
\type{USE STAFF}.

\ATEHead{The Functions}
\texttt{\$} is true when one string occurs in another; \fn{AT()} gives the place
where it begins; \fn{STUFF()} replaces part of a string with another:
\DotExample{examples/p17-strings}
The first letter of a name, in capitals, joined to the rest in small letters,
gives the name as a letter prints it.

\ATEHead{Salutations}
\file{NAMEFMT.PRG} prints the salutation of a letter to each employee:
\ProgramListing{NAMEFMT.PRG}
\begin{Steps}
\item Type \type{DO NAMEFMT}.
\end{Steps}
\DotExample{examples/p17-namefmt}
\fn{TRIM()} removes the blanks at the end of each field, so that the first and
the last name are separated by one blank and the comma follows the name.

\ATEHead{Where Do You Go From Here?}
Chapter~18 changes many records in one program.

\ATEHead{Summary}
\begin{enumerate}
\item \fn{LEFT()} and \fn{SUBSTR()} take strings apart; \texttt{+} joins them.
\item \fn{UPPER()} and \fn{LOWER()} change case; \fn{TRIM()} removes trailing
  blanks.
\item \texttt{\$}, \fn{AT()} and \fn{STUFF()} search and replace.
\end{enumerate}
